% =====================================================================
%  preamble.tex  —  shared preamble for Part I (physics.tex) and
%  Part II (math.tex) of "Bell Correlation as a Foundational Principle".
%
%  RECONSTRUCTED minimal preamble.  The original was referenced in
%  README_build_notes.md but never created as a file; physics/math were
%  only ever compiled against a throwaway stub.  This file collects the
%  macros and environments that physics.tex and math.tex are known to use
%  (from project records): a single shared theorem-numbering scheme across
%  both parts, the status labels, the NSA symbols (st, ord, *R), and the
%  categorical chain C_corr -> C_dist -> C_act -> C_eq.
%
%  BUILD (per README_build_notes.md), bidirectional xr cross-references:
%    each of physics.tex / math.tex does % =====================================================================
%  preamble.tex  —  shared preamble for Part I (physics.tex) and
%  Part II (math.tex) of "Bell Correlation as a Foundational Principle".
%
%  RECONSTRUCTED minimal preamble.  The original was referenced in
%  README_build_notes.md but never created as a file; physics/math were
%  only ever compiled against a throwaway stub.  This file collects the
%  macros and environments that physics.tex and math.tex are known to use
%  (from project records): a single shared theorem-numbering scheme across
%  both parts, the status labels, the NSA symbols (st, ord, *R), and the
%  categorical chain C_corr -> C_dist -> C_act -> C_eq.
%
%  BUILD (per README_build_notes.md), bidirectional xr cross-references:
%    each of physics.tex / math.tex does % =====================================================================
%  preamble.tex  —  shared preamble for Part I (physics.tex) and
%  Part II (math.tex) of "Bell Correlation as a Foundational Principle".
%
%  RECONSTRUCTED minimal preamble.  The original was referenced in
%  README_build_notes.md but never created as a file; physics/math were
%  only ever compiled against a throwaway stub.  This file collects the
%  macros and environments that physics.tex and math.tex are known to use
%  (from project records): a single shared theorem-numbering scheme across
%  both parts, the status labels, the NSA symbols (st, ord, *R), and the
%  categorical chain C_corr -> C_dist -> C_act -> C_eq.
%
%  BUILD (per README_build_notes.md), bidirectional xr cross-references:
%    each of physics.tex / math.tex does % =====================================================================
%  preamble.tex  —  shared preamble for Part I (physics.tex) and
%  Part II (math.tex) of "Bell Correlation as a Foundational Principle".
%
%  RECONSTRUCTED minimal preamble.  The original was referenced in
%  README_build_notes.md but never created as a file; physics/math were
%  only ever compiled against a throwaway stub.  This file collects the
%  macros and environments that physics.tex and math.tex are known to use
%  (from project records): a single shared theorem-numbering scheme across
%  both parts, the status labels, the NSA symbols (st, ord, *R), and the
%  categorical chain C_corr -> C_dist -> C_act -> C_eq.
%
%  BUILD (per README_build_notes.md), bidirectional xr cross-references:
%    each of physics.tex / math.tex does \input{preamble} then
%    \usepackage{xr}\externaldocument[I-]{physics}  (in math.tex)
%    \usepackage{xr}\externaldocument[II-]{math}     (in physics.tex)
%    Compile order: bootstrap each once (pdflatex) to make .aux, then
%    latexmk both twice to converge cross-document refs.
%
%  This preamble is deliberately conservative: it defines every macro the
%  two documents are known to reference, plus the common analysis symbols,
%  so a missing \input does not break the build.  If physics/math use a
%  macro not defined here, add it here (this is the single shared source).
%  files own \documentclass (verified: both real files issue
%  \documentclass[11pt]{article} on line 1), so this preamble must NOT.
% =====================================================================

% ---- encoding / fonts / layout ----
\usepackage[T1]{fontenc}
\usepackage[utf8]{inputenc}
\usepackage[margin=1.1in]{geometry}

% ---- math ----
\usepackage{amsmath,amsthm,amssymb,mathtools}
\usepackage{bm}

% ---- graphics / tables / misc ----
\usepackage{graphicx}
\usepackage{booktabs}
\usepackage{enumitem}
\usepackage{xcolor}

% ---- hyperref last (before xr in the document) ----
\usepackage[colorlinks=true,linkcolor=blue,citecolor=blue,urlcolor=blue]{hyperref}

% =====================================================================
%  Theorem environments — ONE shared numbering scheme across both parts.
%  All share the counter "thmcounter" numbered within section, so Part I
%  and Part II do not collide when cross-referenced.
% =====================================================================
\newtheorem{theorem}{Theorem}[section]
\newtheorem{proposition}[theorem]{Proposition}
\newtheorem{lemma}[theorem]{Lemma}
\newtheorem{corollary}[theorem]{Corollary}
\theoremstyle{definition}
\newtheorem{definition}[theorem]{Definition}
\newtheorem{example}[theorem]{Example}
\theoremstyle{remark}
\newtheorem{remark}[theorem]{Remark}
\newtheorem{note}[theorem]{Note}
\newtheorem{principle}[theorem]{Principle}

% =====================================================================
%  Status labels (README §7: TeX holds only established statements;
%  these tag the epistemic status inline).
% =====================================================================
\newcommand{\established}{\textcolor{green!55!black}{\textsf{[established]}}}
\newcommand{\conditional}{\textcolor{orange!85!black}{\textsf{[conditional]}}}
\newcommand{\deferred}{\textcolor{blue!70!black}{\textsf{[deferred]}}}
\newcommand{\TODO}{\textcolor{red!80!black}{\textsf{[TODO]}}}

% =====================================================================
%  Nonstandard-analysis symbols.
%  *R = hyperreals; st = standard-part map; ord = infinitesimal order;
%  fin = finite hyperreals. eps is the fixed positive infinitesimal
%  lattice spacing (with eps^2 > 0 still infinitesimal).
% =====================================================================
\newcommand{\stR}{{}^{*}\mathbb{R}}        % hyperreals
\newcommand{\stN}{{}^{*}\mathbb{N}}        % hypernaturals
\DeclareMathOperator{\st}{st}              % standard-part map
\DeclareMathOperator{\ord}{ord}            % infinitesimal order
\DeclareMathOperator{\fin}{fin}            % finite hyperreals
\newcommand{\eps}{\varepsilon}

% measurement map  M(f) = st( f / eps^{ord f} )
\newcommand{\Meas}{\mathcal{M}}

% =====================================================================
%  Categorical chain  C_corr -> C_dist -> C_act -> C_eq.
% =====================================================================
\newcommand{\Ccorr}{\mathcal{C}_{\mathrm{corr}}}
\newcommand{\Cdist}{\mathcal{C}_{\mathrm{dist}}}
\newcommand{\Cact}{\mathcal{C}_{\mathrm{act}}}
\newcommand{\Ceq}{\mathcal{C}_{\mathrm{eq}}}

% =====================================================================
%  Common shorthands used across the two parts.
% =====================================================================
\newcommand{\R}{\mathbb{R}}
\newcommand{\N}{\mathbb{N}}
\newcommand{\Z}{\mathbb{Z}}
\newcommand{\C}{\mathbb{C}}
\newcommand{\Q}{\mathbb{Q}}
\newcommand{\corr}{C}                      % correlation C(theta) = -cos theta
\newcommand{\dist}{D}                      % relational distance D = -log|C|
\newcommand{\elliptope}{\mathcal{E}}
\newcommand{\Gram}{G}
\newcommand{\inner}[2]{\langle #1, #2 \rangle}
\newcommand{\abs}[1]{\left| #1 \right|}
\newcommand{\norm}[1]{\left\lVert #1 \right\rVert}
\newcommand{\set}[1]{\left\{ #1 \right\}}
\DeclareMathOperator{\rank}{rank}
\DeclareMathOperator{\tr}{tr}
\DeclareMathOperator{\diag}{diag}
\DeclareMathOperator*{\argmin}{arg\,min}
\DeclareMathOperator*{\argmax}{arg\,max}

% Bell-framework specifics: half-angle delta, Boole quantities F_eps, the
% window depth m = min_eps F_eps.
\newcommand{\Fbool}{F}
\newcommand{\windowdepth}{m}

% =====================================================================
%  Macros required by the real physics.tex / math.tex (verified by macro
%  extraction). Meanings inferred from usage context.
% =====================================================================
\newcommand{\RR}{\mathbb{R}}               % reals (used as \RR)
\newcommand{\ZZ}{\mathbb{Z}}               % integers (Z = pi_1(S^1))
\newcommand{\Cmat}{C}                       % correlation matrix / value, used as |\Cmat|
\newcommand{\dent}{d}                       % relational distance d=-log|C|
\newcommand{\Tcirc}{\mathbb{T}}            % torus; used as \Tcirc^{3}, \Tcirc^{p}
\newcommand{\Scirc}{\mathbb{S}}            % circle S^1; Z \cong pi_1(\Scirc)
\newcommand{\LamQCD}{\Lambda_{\mathrm{QCD}}}% QCD scale
\newcommand{\Eang}{E}                       % E_ang: correlation E(theta) as a function of angle; \Eang(theta)=cos 2theta/(2N)

% end of preamble.tex
 then
%    \usepackage{xr}\externaldocument[I-]{physics}  (in math.tex)
%    \usepackage{xr}\externaldocument[II-]{math}     (in physics.tex)
%    Compile order: bootstrap each once (pdflatex) to make .aux, then
%    latexmk both twice to converge cross-document refs.
%
%  This preamble is deliberately conservative: it defines every macro the
%  two documents are known to reference, plus the common analysis symbols,
%  so a missing \input does not break the build.  If physics/math use a
%  macro not defined here, add it here (this is the single shared source).
%  files own \documentclass (verified: both real files issue
%  \documentclass[11pt]{article} on line 1), so this preamble must NOT.
% =====================================================================

% ---- encoding / fonts / layout ----
\usepackage[T1]{fontenc}
\usepackage[utf8]{inputenc}
\usepackage[margin=1.1in]{geometry}

% ---- math ----
\usepackage{amsmath,amsthm,amssymb,mathtools}
\usepackage{bm}

% ---- graphics / tables / misc ----
\usepackage{graphicx}
\usepackage{booktabs}
\usepackage{enumitem}
\usepackage{xcolor}

% ---- hyperref last (before xr in the document) ----
\usepackage[colorlinks=true,linkcolor=blue,citecolor=blue,urlcolor=blue]{hyperref}

% =====================================================================
%  Theorem environments — ONE shared numbering scheme across both parts.
%  All share the counter "thmcounter" numbered within section, so Part I
%  and Part II do not collide when cross-referenced.
% =====================================================================
\newtheorem{theorem}{Theorem}[section]
\newtheorem{proposition}[theorem]{Proposition}
\newtheorem{lemma}[theorem]{Lemma}
\newtheorem{corollary}[theorem]{Corollary}
\theoremstyle{definition}
\newtheorem{definition}[theorem]{Definition}
\newtheorem{example}[theorem]{Example}
\theoremstyle{remark}
\newtheorem{remark}[theorem]{Remark}
\newtheorem{note}[theorem]{Note}
\newtheorem{principle}[theorem]{Principle}

% =====================================================================
%  Status labels (README §7: TeX holds only established statements;
%  these tag the epistemic status inline).
% =====================================================================
\newcommand{\established}{\textcolor{green!55!black}{\textsf{[established]}}}
\newcommand{\conditional}{\textcolor{orange!85!black}{\textsf{[conditional]}}}
\newcommand{\deferred}{\textcolor{blue!70!black}{\textsf{[deferred]}}}
\newcommand{\TODO}{\textcolor{red!80!black}{\textsf{[TODO]}}}

% =====================================================================
%  Nonstandard-analysis symbols.
%  *R = hyperreals; st = standard-part map; ord = infinitesimal order;
%  fin = finite hyperreals. eps is the fixed positive infinitesimal
%  lattice spacing (with eps^2 > 0 still infinitesimal).
% =====================================================================
\newcommand{\stR}{{}^{*}\mathbb{R}}        % hyperreals
\newcommand{\stN}{{}^{*}\mathbb{N}}        % hypernaturals
\DeclareMathOperator{\st}{st}              % standard-part map
\DeclareMathOperator{\ord}{ord}            % infinitesimal order
\DeclareMathOperator{\fin}{fin}            % finite hyperreals
\newcommand{\eps}{\varepsilon}

% measurement map  M(f) = st( f / eps^{ord f} )
\newcommand{\Meas}{\mathcal{M}}

% =====================================================================
%  Categorical chain  C_corr -> C_dist -> C_act -> C_eq.
% =====================================================================
\newcommand{\Ccorr}{\mathcal{C}_{\mathrm{corr}}}
\newcommand{\Cdist}{\mathcal{C}_{\mathrm{dist}}}
\newcommand{\Cact}{\mathcal{C}_{\mathrm{act}}}
\newcommand{\Ceq}{\mathcal{C}_{\mathrm{eq}}}

% =====================================================================
%  Common shorthands used across the two parts.
% =====================================================================
\newcommand{\R}{\mathbb{R}}
\newcommand{\N}{\mathbb{N}}
\newcommand{\Z}{\mathbb{Z}}
\newcommand{\C}{\mathbb{C}}
\newcommand{\Q}{\mathbb{Q}}
\newcommand{\corr}{C}                      % correlation C(theta) = -cos theta
\newcommand{\dist}{D}                      % relational distance D = -log|C|
\newcommand{\elliptope}{\mathcal{E}}
\newcommand{\Gram}{G}
\newcommand{\inner}[2]{\langle #1, #2 \rangle}
\newcommand{\abs}[1]{\left| #1 \right|}
\newcommand{\norm}[1]{\left\lVert #1 \right\rVert}
\newcommand{\set}[1]{\left\{ #1 \right\}}
\DeclareMathOperator{\rank}{rank}
\DeclareMathOperator{\tr}{tr}
\DeclareMathOperator{\diag}{diag}
\DeclareMathOperator*{\argmin}{arg\,min}
\DeclareMathOperator*{\argmax}{arg\,max}

% Bell-framework specifics: half-angle delta, Boole quantities F_eps, the
% window depth m = min_eps F_eps.
\newcommand{\Fbool}{F}
\newcommand{\windowdepth}{m}

% =====================================================================
%  Macros required by the real physics.tex / math.tex (verified by macro
%  extraction). Meanings inferred from usage context.
% =====================================================================
\newcommand{\RR}{\mathbb{R}}               % reals (used as \RR)
\newcommand{\ZZ}{\mathbb{Z}}               % integers (Z = pi_1(S^1))
\newcommand{\Cmat}{C}                       % correlation matrix / value, used as |\Cmat|
\newcommand{\dent}{d}                       % relational distance d=-log|C|
\newcommand{\Tcirc}{\mathbb{T}}            % torus; used as \Tcirc^{3}, \Tcirc^{p}
\newcommand{\Scirc}{\mathbb{S}}            % circle S^1; Z \cong pi_1(\Scirc)
\newcommand{\LamQCD}{\Lambda_{\mathrm{QCD}}}% QCD scale
\newcommand{\Eang}{E}                       % E_ang: correlation E(theta) as a function of angle; \Eang(theta)=cos 2theta/(2N)

% end of preamble.tex
 then
%    \usepackage{xr}\externaldocument[I-]{physics}  (in math.tex)
%    \usepackage{xr}\externaldocument[II-]{math}     (in physics.tex)
%    Compile order: bootstrap each once (pdflatex) to make .aux, then
%    latexmk both twice to converge cross-document refs.
%
%  This preamble is deliberately conservative: it defines every macro the
%  two documents are known to reference, plus the common analysis symbols,
%  so a missing \input does not break the build.  If physics/math use a
%  macro not defined here, add it here (this is the single shared source).
%  files own \documentclass (verified: both real files issue
%  \documentclass[11pt]{article} on line 1), so this preamble must NOT.
% =====================================================================

% ---- encoding / fonts / layout ----
\usepackage[T1]{fontenc}
\usepackage[utf8]{inputenc}
\usepackage[margin=1.1in]{geometry}

% ---- math ----
\usepackage{amsmath,amsthm,amssymb,mathtools}
\usepackage{bm}

% ---- graphics / tables / misc ----
\usepackage{graphicx}
\usepackage{booktabs}
\usepackage{enumitem}
\usepackage{xcolor}

% ---- hyperref last (before xr in the document) ----
\usepackage[colorlinks=true,linkcolor=blue,citecolor=blue,urlcolor=blue]{hyperref}

% =====================================================================
%  Theorem environments — ONE shared numbering scheme across both parts.
%  All share the counter "thmcounter" numbered within section, so Part I
%  and Part II do not collide when cross-referenced.
% =====================================================================
\newtheorem{theorem}{Theorem}[section]
\newtheorem{proposition}[theorem]{Proposition}
\newtheorem{lemma}[theorem]{Lemma}
\newtheorem{corollary}[theorem]{Corollary}
\theoremstyle{definition}
\newtheorem{definition}[theorem]{Definition}
\newtheorem{example}[theorem]{Example}
\theoremstyle{remark}
\newtheorem{remark}[theorem]{Remark}
\newtheorem{note}[theorem]{Note}
\newtheorem{principle}[theorem]{Principle}

% =====================================================================
%  Status labels (README §7: TeX holds only established statements;
%  these tag the epistemic status inline).
% =====================================================================
\newcommand{\established}{\textcolor{green!55!black}{\textsf{[established]}}}
\newcommand{\conditional}{\textcolor{orange!85!black}{\textsf{[conditional]}}}
\newcommand{\deferred}{\textcolor{blue!70!black}{\textsf{[deferred]}}}
\newcommand{\TODO}{\textcolor{red!80!black}{\textsf{[TODO]}}}

% =====================================================================
%  Nonstandard-analysis symbols.
%  *R = hyperreals; st = standard-part map; ord = infinitesimal order;
%  fin = finite hyperreals. eps is the fixed positive infinitesimal
%  lattice spacing (with eps^2 > 0 still infinitesimal).
% =====================================================================
\newcommand{\stR}{{}^{*}\mathbb{R}}        % hyperreals
\newcommand{\stN}{{}^{*}\mathbb{N}}        % hypernaturals
\DeclareMathOperator{\st}{st}              % standard-part map
\DeclareMathOperator{\ord}{ord}            % infinitesimal order
\DeclareMathOperator{\fin}{fin}            % finite hyperreals
\newcommand{\eps}{\varepsilon}

% measurement map  M(f) = st( f / eps^{ord f} )
\newcommand{\Meas}{\mathcal{M}}

% =====================================================================
%  Categorical chain  C_corr -> C_dist -> C_act -> C_eq.
% =====================================================================
\newcommand{\Ccorr}{\mathcal{C}_{\mathrm{corr}}}
\newcommand{\Cdist}{\mathcal{C}_{\mathrm{dist}}}
\newcommand{\Cact}{\mathcal{C}_{\mathrm{act}}}
\newcommand{\Ceq}{\mathcal{C}_{\mathrm{eq}}}

% =====================================================================
%  Common shorthands used across the two parts.
% =====================================================================
\newcommand{\R}{\mathbb{R}}
\newcommand{\N}{\mathbb{N}}
\newcommand{\Z}{\mathbb{Z}}
\newcommand{\C}{\mathbb{C}}
\newcommand{\Q}{\mathbb{Q}}
\newcommand{\corr}{C}                      % correlation C(theta) = -cos theta
\newcommand{\dist}{D}                      % relational distance D = -log|C|
\newcommand{\elliptope}{\mathcal{E}}
\newcommand{\Gram}{G}
\newcommand{\inner}[2]{\langle #1, #2 \rangle}
\newcommand{\abs}[1]{\left| #1 \right|}
\newcommand{\norm}[1]{\left\lVert #1 \right\rVert}
\newcommand{\set}[1]{\left\{ #1 \right\}}
\DeclareMathOperator{\rank}{rank}
\DeclareMathOperator{\tr}{tr}
\DeclareMathOperator{\diag}{diag}
\DeclareMathOperator*{\argmin}{arg\,min}
\DeclareMathOperator*{\argmax}{arg\,max}

% Bell-framework specifics: half-angle delta, Boole quantities F_eps, the
% window depth m = min_eps F_eps.
\newcommand{\Fbool}{F}
\newcommand{\windowdepth}{m}

% =====================================================================
%  Macros required by the real physics.tex / math.tex (verified by macro
%  extraction). Meanings inferred from usage context.
% =====================================================================
\newcommand{\RR}{\mathbb{R}}               % reals (used as \RR)
\newcommand{\ZZ}{\mathbb{Z}}               % integers (Z = pi_1(S^1))
\newcommand{\Cmat}{C}                       % correlation matrix / value, used as |\Cmat|
\newcommand{\dent}{d}                       % relational distance d=-log|C|
\newcommand{\Tcirc}{\mathbb{T}}            % torus; used as \Tcirc^{3}, \Tcirc^{p}
\newcommand{\Scirc}{\mathbb{S}}            % circle S^1; Z \cong pi_1(\Scirc)
\newcommand{\LamQCD}{\Lambda_{\mathrm{QCD}}}% QCD scale
\newcommand{\Eang}{E}                       % E_ang: correlation E(theta) as a function of angle; \Eang(theta)=cos 2theta/(2N)

% end of preamble.tex
 then
%    \usepackage{xr}\externaldocument[I-]{physics}  (in math.tex)
%    \usepackage{xr}\externaldocument[II-]{math}     (in physics.tex)
%    Compile order: bootstrap each once (pdflatex) to make .aux, then
%    latexmk both twice to converge cross-document refs.
%
%  This preamble is deliberately conservative: it defines every macro the
%  two documents are known to reference, plus the common analysis symbols,
%  so a missing \input does not break the build.  If physics/math use a
%  macro not defined here, add it here (this is the single shared source).
%  files own \documentclass (verified: both real files issue
%  \documentclass[11pt]{article} on line 1), so this preamble must NOT.
% =====================================================================

% ---- encoding / fonts / layout ----
\usepackage[T1]{fontenc}
\usepackage[utf8]{inputenc}
\usepackage[margin=1.1in]{geometry}

% ---- math ----
\usepackage{amsmath,amsthm,amssymb,mathtools}
\usepackage{bm}

% ---- graphics / tables / misc ----
\usepackage{graphicx}
\usepackage{booktabs}
\usepackage{enumitem}
\usepackage{xcolor}

% ---- hyperref last (before xr in the document) ----
\usepackage[colorlinks=true,linkcolor=blue,citecolor=blue,urlcolor=blue]{hyperref}

% =====================================================================
%  Theorem environments — ONE shared numbering scheme across both parts.
%  All share the counter "thmcounter" numbered within section, so Part I
%  and Part II do not collide when cross-referenced.
% =====================================================================
\newtheorem{theorem}{Theorem}[section]
\newtheorem{proposition}[theorem]{Proposition}
\newtheorem{lemma}[theorem]{Lemma}
\newtheorem{corollary}[theorem]{Corollary}
\theoremstyle{definition}
\newtheorem{definition}[theorem]{Definition}
\newtheorem{example}[theorem]{Example}
\theoremstyle{remark}
\newtheorem{remark}[theorem]{Remark}
\newtheorem{note}[theorem]{Note}
\newtheorem{principle}[theorem]{Principle}

% =====================================================================
%  Status labels (README §7: TeX holds only established statements;
%  these tag the epistemic status inline).
% =====================================================================
\newcommand{\established}{\textcolor{green!55!black}{\textsf{[established]}}}
\newcommand{\conditional}{\textcolor{orange!85!black}{\textsf{[conditional]}}}
\newcommand{\deferred}{\textcolor{blue!70!black}{\textsf{[deferred]}}}
\newcommand{\TODO}{\textcolor{red!80!black}{\textsf{[TODO]}}}

% =====================================================================
%  Nonstandard-analysis symbols.
%  *R = hyperreals; st = standard-part map; ord = infinitesimal order;
%  fin = finite hyperreals. eps is the fixed positive infinitesimal
%  lattice spacing (with eps^2 > 0 still infinitesimal).
% =====================================================================
\newcommand{\stR}{{}^{*}\mathbb{R}}        % hyperreals
\newcommand{\stN}{{}^{*}\mathbb{N}}        % hypernaturals
\DeclareMathOperator{\st}{st}              % standard-part map
\DeclareMathOperator{\ord}{ord}            % infinitesimal order
\DeclareMathOperator{\fin}{fin}            % finite hyperreals
\newcommand{\eps}{\varepsilon}

% measurement map  M(f) = st( f / eps^{ord f} )
\newcommand{\Meas}{\mathcal{M}}

% =====================================================================
%  Categorical chain  C_corr -> C_dist -> C_act -> C_eq.
% =====================================================================
\newcommand{\Ccorr}{\mathcal{C}_{\mathrm{corr}}}
\newcommand{\Cdist}{\mathcal{C}_{\mathrm{dist}}}
\newcommand{\Cact}{\mathcal{C}_{\mathrm{act}}}
\newcommand{\Ceq}{\mathcal{C}_{\mathrm{eq}}}

% =====================================================================
%  Common shorthands used across the two parts.
% =====================================================================
\newcommand{\R}{\mathbb{R}}
\newcommand{\N}{\mathbb{N}}
\newcommand{\Z}{\mathbb{Z}}
\newcommand{\C}{\mathbb{C}}
\newcommand{\Q}{\mathbb{Q}}
\newcommand{\corr}{C}                      % correlation C(theta) = -cos theta
\newcommand{\dist}{D}                      % relational distance D = -log|C|
\newcommand{\elliptope}{\mathcal{E}}
\newcommand{\Gram}{G}
\newcommand{\inner}[2]{\langle #1, #2 \rangle}
\newcommand{\abs}[1]{\left| #1 \right|}
\newcommand{\norm}[1]{\left\lVert #1 \right\rVert}
\newcommand{\set}[1]{\left\{ #1 \right\}}
\DeclareMathOperator{\rank}{rank}
\DeclareMathOperator{\tr}{tr}
\DeclareMathOperator{\diag}{diag}
\DeclareMathOperator*{\argmin}{arg\,min}
\DeclareMathOperator*{\argmax}{arg\,max}

% Bell-framework specifics: half-angle delta, Boole quantities F_eps, the
% window depth m = min_eps F_eps.
\newcommand{\Fbool}{F}
\newcommand{\windowdepth}{m}

% =====================================================================
%  Macros required by the real physics.tex / math.tex (verified by macro
%  extraction). Meanings inferred from usage context.
% =====================================================================
\newcommand{\RR}{\mathbb{R}}               % reals (used as \RR)
\newcommand{\ZZ}{\mathbb{Z}}               % integers (Z = pi_1(S^1))
\newcommand{\Cmat}{C}                       % correlation matrix / value, used as |\Cmat|
\newcommand{\dent}{d}                       % relational distance d=-log|C|
\newcommand{\Tcirc}{\mathbb{T}}            % torus; used as \Tcirc^{3}, \Tcirc^{p}
\newcommand{\Scirc}{\mathbb{S}}            % circle S^1; Z \cong pi_1(\Scirc)
\newcommand{\LamQCD}{\Lambda_{\mathrm{QCD}}}% QCD scale
\newcommand{\Eang}{E}                       % E_ang: correlation E(theta) as a function of angle; \Eang(theta)=cos 2theta/(2N)

% end of preamble.tex
